\section{Limitations}
\label{sec:limitations}

\begin{enumerate}
  \item \exactmatch's own audited $\alpha,\beta$ do not yet exist---Analysis
  C is the single largest open piece of this paper, \pending{458 rows,
  preregistered and in progress}.

  \item The correlated-error cancellation result (Section~\ref{sec:theory})
  is verified by simulation against its own closed form, not fit to any
  real $(\pi, p_{\text{hard}}, p_{\text{easy}})$---no dataset has produced
  enough scorer-wrong events to fit it (0/69 real T-C events).

  \item The \scoreboxed-path comparator's behavior was verified correct on
  Linux via continuous integration, resolving the original Windows-only bug
  finding, but two smaller, named structural limitations found in the
  adversarial suite (interval/set-notation reordering; unit-word stripping)
  remain unfixed---flagged, not silently left stale.

  \item The scorer-choice generalization claim
  (Section~\ref{sec:results-generalization}) is deliberately narrow: one
  benchmark in the strong sense. IFEval offers a second, real data point
  (2.60\% disagreement between its \code{strict\_acc}/\code{loose\_acc}
  settings) but is structurally different---those are two leniency settings
  of the \emph{same} verifier, not two independently-implemented libraries,
  so it corroborates without replicating. A second benchmark with two
  genuinely independent scoring libraries was searched for and not found
  within this paper's currently-accessible data sources---the ``evaluation
  pipelines as imperfect measurement systems'' claim is accordingly stated
  as established for MATH-Hard and suggestive for IFEval, not as a general
  property of automated evaluation.

  \item The asymptotic coverage properties of the bootstrap-based
  flip-budget interval (Sections~\ref{sec:theory},~\ref{sec:results-flip-budget})
  are not proven for this specific nonlinear, potentially-empty-set
  estimand---a real, named gap for a statistics-focused venue, not claimed
  resolved.

  \item Naive-Neyman-allocation-fails is demonstrated on this paper's own
  real, sparse data; whether it generalizes to audits of other benchmarks
  with different sparsity patterns is untested.

  \item The T-C round's 69 additional audited items
  (Section~\ref{sec:empirical-setup}) were analyzed for their standalone
  effect on pooling but never merged into the canonical per-model pipeline;
  every headline number in Section~\ref{sec:results} reflects only the
  original 350-usable-item audit. Merging is mechanical (re-run the pooling
  with both label sets combined and regenerate the downstream chain) but
  has not been done---named explicitly rather than left as a silent
  staleness gap.

  \item \textbf{The underlying audits are thin, and the identified-set width
  this paper reports is partly a statement about that thinness, not only
  about true scorer error.} The median per-model audit size is small
  ($n=11$ in the real-data negative-control check,
  Section~\ref{sec:audit-design}); 15 of 17 qualifying models have exactly
  zero observed false-credit events. The negative-controls exercise shows a
  \emph{perfect} scorer audited at this paper's realistic sample sizes
  still produces a large dominance ratio---the headline finding is real,
  but a reader should not infer ``this benchmark's scorer is unusually
  error-prone'' from it alone; the correct inference is ``this benchmark's
  scorer error has not been pinned down as tightly as its sampling
  variability has,'' which is this paper's actual claim.

  \item \textbf{This paper depends on the assumption that a human reader
  can, in principle, determine ground-truth correctness for a MATH-Hard
  item}---the $(\alpha,\beta)$ framework bounds \emph{scorer} error against
  that reference, not against some deeper notion of mathematical truth
  independent of any reader. T-D's sensitivity bound
  (Section~\ref{sec:theory}) addresses reader error explicitly and finds no
  concerning breakdown point at realistic human-error rates, but the
  framework as a whole is not agnostic to this assumption; it is a standard
  one in the measurement-error literature and is stated here rather than
  left implicit.

  \item \textbf{Several of this paper's supporting findings are bug
  reports, not methodological contributions, and should be read as such.}
  The Windows-specific comparator failure, the empty-identified-set
  silent-inversion bug (Section~\ref{sec:problem-formulation}), the
  Horvitz--Thompson pooling bias in the audit-expansion pipeline, and a
  third-party library's own multiprocessing-timeout fragility are all real,
  disclosed, and load-bearing for this paper's own internal correctness---but
  none of them is a novelty claim. They motivate and validate the broader
  argument (that evaluation pipelines are imperfect measurement systems
  worth auditing) without being examples of it in the formal sense;
  conflating ``we found and fixed $N$ bugs in our own tooling'' with ``we
  contributed $N$ results'' would overstate this paper's contribution, and
  it does not make that claim (Section~\ref{sec:theory},
  Appendix~\ref{sec:appendix-novelty}).
\end{enumerate}
